\section*{Hoofdstuk \ref{ch:b2:deep-time-climate} --- Het klimaat over de geologische tijd}

\begin{solution}{exo:b2:deep-time-climate:1}
(a) $900 + 37100 + 700 + 3 = \qty{38703}{PgC}$, ongeveer 65 keer de \qty{591}{PgC} van de lucht.
(b) $\num{3.5e6}\times0.012011 = \qty{42000}{PgC}$. (c) $591/42000 = \qty{1.4}{\percent}$.
\end{solution}

\begin{solution}{exo:b2:deep-time-climate:2}
(a) Instroom $11.3 + 3 + 12.3 + 8.1 + 2.0 = \qty{36.7}{Tmol.yr^{-1}}$; uitstroom $26 + 16 =
\qty{42}{Tmol.yr^{-1}}$, gelijk binnen hun onzekerheden. (b) $(11.3 + 3)/36.7 = 0.39$: ruim
onder de helft; het grootste deel van de instroom is oude sedimentaire koolstof. (c) $\num{3.5e6}/42 = \qty{8.3e4}{yr}$.
\end{solution}

\begin{solution}{exo:b2:deep-time-climate:3}
(a) \ce{Mg2SiO4 + 4 CO2 + 2 H2O -> 2 Mg^{2+} + 4 HCO3^- + SiO2}. (b) Met \ce{2 Mg^{2+} + 4 HCO3^-
-> 2 MgCO3 + 2 CO2 + 2 H2O} is de nettoreactie \ce{Mg2SiO4 + 2 CO2 -> 2 MgCO3 + SiO2}: twee mol
koolstofdioxide per mol forsteriet. (c) $\qty{1}{PgC} = 10^{15}/12.011 =
\qty{8.33e13}{mol}$ koolstof, waarvoor \qty{4.16e13}{mol} forsteriet nodig is, $4.16\times10^{13}
\times\qty{140.7}{g} = \qty{5.86}{Pg}$, bijna zes miljard ton.
\end{solution}

\begin{solution}{exo:b2:deep-time-climate:4}
Hirnantiaanse vergletsjering, 445 tot \qty{443}{Ma}, ijskast; grootste ijsuitbreiding van het late Paleozoïcum, 307 tot
\qty{299}{Ma}, ijskast; Turonien, 94 tot \qty{89}{Ma}, broeikas (de warmste van het Fanerozoïcum);
PETM, 55.9 tot \qty{55.7}{Ma}, broeikas; EECO, 53 tot \qty{49}{Ma}, broeikas; ontstaan van de
Antarctische ijskap, ongeveer \qty{34}{Ma}, begin van de huidige ijskast; noordelijke ijskappen,
\qty{2.7}{Ma}, ijskast.
\end{solution}

\begin{solution}{exo:b2:deep-time-climate:5}
(a) $63000/(8.314\times288^2) = \qty{0.0914}{K^{-1}}$, wat afgerond 0.09 geeft. (b) $74000/(8.314
\times288^2) = \qty{0.107}{K^{-1}}$. (c) $\ln1.2/0.09 = \qty{2.03}{K}$ en $\ln 1.2/0.107 =
\qty{1.70}{K}$. (d) In de stationaire toestand moet de verwering gelijk zijn aan de ontgassing, wat $e^{k
\Delta T}$ en dus $\Delta T$ vastlegt; de gevoeligheid bepaalt alleen hoeveel koolstofdioxide die
opwarming geeft.
\end{solution}

\begin{solution}{exo:b2:deep-time-climate:6}
Hier is $k_T = 0.0914$. (a) Exact $\exp[7577.5\times(1/288 - 1/293)] = 1.567$; benadering
$e^{0.0914\times5} = 1.579$: \qty{0.8}{\percent} te hoog. (b) Exact $e^{2.482} = 11.97$;
benadering $e^{2.741} = 15.50$: \qty{30}{\percent} te hoog. (c) $\ln(\text{verhouding}) = (E_{\text
a}/R)[x/T_0^2 - x/(T_0(T_0 + x))] = (E_{\text a}/R)\,x^2/(T_0^2(T_0 + x))$. Stel dit gelijk aan $\ln1.1$:
$x^2 = 1.043\,(288 + x)$, dus $x = \qty{17.9}{K}$. Voorbij ongeveer \qty{18}{K} opwarming
overschat de exponentiële vorm de verwering met meer dan een tiende.
\end{solution}

\begin{solution}{exo:b2:deep-time-climate:7}
(a) $1444/276 = 5.23$, en $3.93\log_2 5.23 = \qty{9.4}{W.m^{-2}}$. (b) Bijna precies de
\qty{9}{W.m^{-2}} van de zon: de twee forceringen heffen elkaar op. (c) $(1444 - 276)/419.5 =
\qty{2.8}{ppm}$ per miljoen jaar. De \qty{3.4}{ppm} komt van een rechte die aan alle
gegevens is aangepast, die elke schatting laat meewegen in plaats van de twee uiteinden van de afgevlakte kromme.
\end{solution}

\begin{solution}{exo:b2:deep-time-climate:8}
(a) Van $3000/20000 = \qty{0.15}{PgC.yr^{-1}}$ tot $7000/3000 = \qty{2.3}{PgC.yr^{-1}}$, een bereik
dat de 0.3 tot 1.5 van het IPCC omvat. (b) $35\times12.011/44.01 = \qty{9.6}{PgC.yr^{-1}}$, 6 tot
32 keer de 0.3 tot \qty{1.5}{PgC.yr^{-1}} van het PETM. (c) $3.93\log_2(2000/900) = \qty{4.5}{W.m^{-2}}$,
en $3\log_2 2.22 = \qty{3.5}{\celsius}$: ongeveer zeventig procent van de \qty{5}{\celsius}. Andere
terugkoppelingen of een andere bron van opwarming, zoals methaan, moeten de rest leveren.
\end{solution}

\begin{solution}{exo:b2:deep-time-climate:9}
(a) Nu is $1.2\,W_0e^{k\Delta T} = V = W_0$, dus $\Delta T = -\ln 1.2/0.09 = \qty{-2.03}{K}$.
(b) $278\times2^{-2.03/3} = \qty{174}{ppm}$. (c) Nee: in elke stationaire toestand is de verwering gelijk aan de
ontgassing, die niet is veranderd. De hogere verweerbaarheid wordt gecompenseerd door een kouder klimaat.
\end{solution}

\begin{solution}{exo:b2:deep-time-climate:10}
(a) $k = 0.119$: $\tau = 591\times0.693/(0.08\times0.119\times3\times0.138) = \qty{104}{kyr}$.
(b) $S = 8$: $\tau = \qty{51.5}{kyr}$. (c) Beide maken de thermostaat sneller, want $\tau \propto
1/(kS)$. $E_{\text a}$ is onzeker tot op een factor van ongeveer twee, $S$ op deze tijdschalen tot op een
vergelijkbare factor (3 tegenover 5 tot \qty{8}{K}): beide wegen ongeveer even zwaar.
\end{solution}

\begin{solution}{exo:b2:deep-time-climate:11}
(a) $0.08\times5000 = \qty{400}{PgC}$ blijft in de lucht: $C/C_0 = 991/591 = 1.677$, of
\qty{466}{ppm}; $\Delta T = 3\log_2 1.677 = \qty{2.24}{K}$. (b) $W - W_0 = 0.138\,(e^{0.09\times2.24}
- 1) = \qty{0.0308}{PgC.yr^{-1}}$, \qty{22}{\percent} meer verwering; $5000/0.0308 =
\qty{162}{kyr}$. (c) $\tau = \qty{137}{kyr}$ en $\tau\ln10 = \qty{315}{kyr}$. (d) De opwarming groeit
slechts met de logaritme van de koolstofdioxide, dus een grote overmaat verhoogt de verwering minder dan de
lineaire wet aanneemt; de verwijdering begint trager.
\end{solution}

\begin{solution}{exo:b2:deep-time-climate:12}
(a) Dicht bij $M^*$ is $\mathrm d\Delta T/\mathrm dM = Sf/(A_0\ln 2\,(1 + fM^*/A_0)) = Sf/(A_0\ln 2\,
C^*/C_0)$ en $\mathrm dW/\mathrm dM = kW\,\mathrm d\Delta T/\mathrm dM$, met daar $W = V$; de
inverse is $\tau'$. (b) $C^*/C_0 = 2^{2.03/3} = 1.597$: $\tau' = 137.4\times1.597/1.2 =
\qty{183}{kyr}$. (c) Trager: bij veel koolstofdioxide warmt elke toegevoegde petagram koolstof minder op,
omdat de forcering logaritmisch is, dus de verwering reageert er minder op.
\end{solution}

\begin{solution}{pb:b2:deep-time-climate:1}
\textbf{1.} $1.3 + 1.3 + 4.5 + 4.2 = \qty{11.3}{Tmol.yr^{-1}}$, $\qty{0.136}{PgC.yr^{-1}}$.
\textbf{2.} $11.5\times0.012011 = \qty{0.138}{PgC.yr^{-1}}$.
\textbf{3.} Ze komen tot op ongeveer \qty{2}{\percent} overeen: de silicaatverwering houdt de vulkanische
ontgassing in evenwicht, zoals de \omterm{def:b2:deep-time-climate:carbon-cycle}{koolstofcyclus op lange termijn} vereist.
\textbf{4.} $\num{5.1e18}/0.02897 = \qty{1.76e20}{mol}$ lucht; $\num{591e15}/12.011 =
\qty{4.92e16}{mol}$ koolstof; $\num{4.92e16}/\num{1.76e20} = \num{2.80e-4}$, \qty{280}{ppm}.
\textbf{5.} $591/0.136 = \qty{4350}{yr}$.
\textbf{6.} $1 + 0.08 = 1.08$: slechts \qty{8}{\percent} meer.
\textbf{7.} $63000/(8.314\times288^2) = \qty{0.0914}{K^{-1}}$, wat afgerond 0.09 geeft.
\textbf{8.} $e^{0.09} = 1.09$: \qty{9}{\percent} meer per kelvin.
\textbf{9.} $\Delta T = S\log_2(1 + fM/A_0)$; de afgeleide ervan bij 0 is $Sf/(A_0\ln2) =
3\times0.08/(591\times0.693) = \qty{5.86e-4}{K.PgC^{-1}}$.
\textbf{10.} $W_0k\times\num{5.86e-4} = 0.138\times0.09\times\num{5.86e-4} = \qty{7.28e-6}{yr^{-1}}$.
\textbf{11.} De inverse ervan is $A_0\ln2/(fkSW_0)$, de $\tau$ van de propositie, want de
gelineariseerde vergelijking is $\mathrm dM/\mathrm dt = -(\mathrm dW/\mathrm dM)\,M$.
\textbf{12.} Alle vier: $\tau$ is omgekeerd evenredig met elk ervan.
\textbf{13.} $1/\num{7.28e-6} = \qty{1.37e5}{yr}$.
\textbf{14.} $137\times0.09/0.119 = \qty{104}{kyr}$.
\textbf{15.} $k_T = 0.107$: $137\times0.09/0.107 = \qty{115}{kyr}$.
\textbf{16.} Korter met een factor 1.5 tot 7: het eendoosmodel en zijn constante $f$ maken de
thermostaat te snel.
\textbf{17.} $\num{50e6}/\num{1.37e5} = 364$: de thermostaat had in elk stadium de tijd om de cyclus in evenwicht te brengen,
dus de daling volgde trage veranderingen van de ontgassing of van de verweerbaarheid.
\textbf{18.} $16\times0.137 = \qty{2.2}{ppm}$.
\textbf{19.} $0.08\times5000 = \qty{400}{PgC}$; $278\times991/591 = \qty{466}{ppm}$.
\textbf{20.} $3\log_2 1.677 = \qty{2.24}{K}$.
\textbf{21.} $0.138\,(e^{0.202} - 1) = \qty{0.0308}{PgC.yr^{-1}}$, \qty{22.3}{\percent}.
\textbf{22.} $137\times\ln10 = \qty{315}{kyr}$.
\textbf{23.} De verwering kan gevoeliger zijn geweest (afvoer, een grotere activeringsenergie, de
trage terugkoppelingen van een grotere aardsysteemgevoeligheid); en organische koolstof op het land en in zee kan
opnieuw zijn gegroeid en sneller koolstof uit de lucht hebben gehaald dan de silicaatverwering.
\textbf{24.} In dit model verwijdert de \omterm{def:b2:deep-time-climate:thermostat}{silicaatverweringsthermostaat} een overmaat koolstof met een
e-foldingtijd van ongeveer \textbf{\qty[text-series-to-math=true]{137}{} duizend jaar}.
\end{solution}
